% Methodology snippet — three paragraphs.
% Numbers from runs/20260519T085446Z_variance_decomposition (Phase 2 decomposition),
% runs/20260519T085744Z_b1_n_ablation (B1 N trajectory),
% runs/20260519T091149Z_paired_vs_unpaired (paired CI comparison).

The eval corpus pools $59$ NASDAQ tickers over the ten months that bracket the February~$2020$ volatility regime change. Cross-sectional return scales differ by more than an order of magnitude between the megacap and small-cap names in the panel, and intraday volatility on a single name can triple between the calm pre-crash window and the four weeks that follow. Each matched-$N$ draw from the eval corpus is therefore a draw from a heterogeneous mixture, and the dispersion of any pooled statistic across resamples reflects the mixture rather than estimator noise alone. Table~\ref{tab:variance-decomposition} quantifies the point. For B1 the pooled CV of $\bar{g}_{rr}$ is $0.65$; restricting both halves of the draw to a single ticker drops it only to $0.64$, and restricting to a single regime drops it to $0.62$. The four columns of the table sit within $0.03$ of each other on B1, B4, and B5. If the high pooled CV were a finite-$N$ artifact, restricting the draw to a much smaller and more homogeneous sub-population would have driven it down sharply; it does not.

Varying $N$ alone forces the same conclusion. Holding the pooled sampler fixed and stepping $N$ through $\{50, 100, 200, 400\}$, the mean of $\bar{g}_{rr}$ for B1 falls from $0.584$ to $0.219$, tracking the $1/\sqrt{N}$ scaling expected of a Wasserstein gap. The CV over the same trajectory is $0.679$, $0.668$, $0.672$, $0.637$. An eightfold increase in $N$ shrinks the CV by less than seven percent. Under any classical estimator-noise model the CV would scale as $1/\sqrt{N}$ and an eightfold $N$ increase would shrink it by a factor close to three. The flatness of the trajectory is direct evidence that the dominant variance source is cross-sectional heterogeneity in the corpus rather than finite-sample noise. The wide marginal CIs that we report on $\bar{g}_{rr}$ are not slack in our estimator; they are the honest width of what real intraday data does across tickers and regimes.

The bootstrap confidence interval on the similarity score $s_b = \bar{g}_{rr} / (\bar{g}_{rr} + \bar{g}_{sr})$ should reflect this heterogeneity but should not be inflated by it gratuitously. The standard unpaired bootstrap draws independent index vectors into $g_{rr}$ and $g_{sr}$, which discards the correlation that the two arms inherit from sharing the real-side index draw at construction. We construct $g_{rr}[i]$ and $g_{sr}[i]$ from the same real-side indices $\mathrm{idx}_a$ and resample a single index vector into both arrays per bootstrap iteration. The point estimate $s_b$ is unchanged. The interval tightens by $10$ to $12$~percent in the cells where the per-$i$ correlation between $g_{rr}$ and $g_{sr}$ is positive and near $0.25$, namely B4 against AIL, B1 against GARCH, and B6 against GARCH. In twelve of the eighteen $(\text{bucket}\times\text{generator})$ cells in the benchmark the paired interval is the narrower one. In the remaining six the per-$i$ correlation is approximately zero or slightly negative; the paired and unpaired intervals are then within five percent of each other in width. The paired estimator is uniformly preferable because it uses information the unpaired estimator throws away.
